\chapter{Application Requirements and Specification}
\label{chap:requirements}

\par This chapter defines the application that constitutes the practical contribution of the thesis. The overall objective is twofold: to deliver a defensible, reproducible deep-learning pipeline trained on the CICIoT2023 dataset, and to expose the resulting model through a live demonstration in which the user can submit captured network traffic and observe per-flow intrusion-detection verdicts. The chapter starts from this objective, identifies the actors and use cases, lists the functional and non-functional requirements that shape the design, and concludes with the system-level specification that constrains the implementation in Chapter~\ref{chap:design_implementation}.

\section{Project Goal and Scope}
\label{sec:project_goal}

\par The deliverable is a real-time intrusion detection system (RT-IDS) for IoT networks built on CICIoT2023, structured as two cooperating subsystems. The training subsystem is a reproducible offline pipeline that ingests the raw CICIoT2023 captures, deduplicates and subsamples them per class, builds three train/validation/test splits, trains a deep neural model at three classification granularities (binary, attack-family, and full multi-class), evaluates it with the metrics expected by the IDS literature, and serialises the small set of artefacts needed for inference. The serving subsystem is a web application that loads those artefacts, accepts either a pre-extracted feature CSV in CIC format or a raw packet capture (in which case a flow-extraction tool derives the feature representation), runs the trained model, and presents per-flow verdicts and an aggregated summary in the browser.

\par The system targets an interactive, single-request setting: the user supplies a sample of traffic, the system classifies it, and the verdict is reviewed on screen. Line-rate gateway protection (sub-millisecond per-packet inference at multi-gigabit throughput) is out of scope. This narrower sense of \enquote{real time} sets the latency thresholds in the non-functional requirements below.

\par All three classification granularities are in scope, because comparing them is itself part of the contribution and matches the way most CICIoT2023 papers report results. The 2-class granularity (Benign versus Attack) is the gate an inline production deployment would actually use. The 8-class granularity (Benign plus the seven attack families: DDoS, DoS, Mirai, Reconnaissance, Spoofing, Web, BruteForce) is the operational view a Security Operations Centre analyst would consume. The 34-class granularity (one label per specific attack variant) is the forensic view used in published comparisons. The same train/validation/test rows are reused across all three granularities, with only the label mapping changing, so the comparison is not contaminated by sampling differences.

\section{Stakeholders and Actors}
\label{sec:actors}

\par The system has two actors. As \emph{researcher}, the author trains and evaluates the models and regenerates the artefacts. As \emph{operator}, the same person (or a reviewer reproducing a result) uploads captured traffic through the web interface and reads the per-flow verdicts. The thesis defence is one instance of the operator role rather than a separate actor. Each role constrains the user-interface design.

\begin{table}[htbp]
\begin{center}
\begin{tabular}{|p{90pt}|p{90pt}|p{170pt}|}
\hline
\textbf{Actor} & \textbf{Role} & \textbf{Primary interaction} \\
\hline
\hline Researcher & Trains and evaluates models, regenerates artefacts. & Runs the training pipeline end-to-end, iterates on splits and granularities, inspects metrics and confusion matrices. \\
\hline Operator & Submits traffic and reviews verdicts (the demonstration and result-verification role). & Launches the web application, selects a model variant, uploads a capture or feature CSV, reads the verdict table and summary, and may reproduce a run to compare against the expected label. \\
\hline
\end{tabular}
\end{center}
\caption{Actors involved in the system and their primary interactions.}
\label{tab:actors}
\end{table}

\par The core classification path is stateless: submitted captures and CSVs are processed in a temporary location and are never persisted (FR-D5). The publicly deployed React frontend layers a thin convenience tier on top of this stateless core (user accounts and a per-user history of past classifications) so that a returning user can review earlier results. That tier stores only \emph{derived result metadata} (the model used, the predicted label, the confidence, and a timestamp), never the uploaded traffic itself. Its design and placement are detailed in Section~\ref{subsec:auth_persistence}.

\section{Functional Requirements}
\label{sec:functional_requirements}

\par The functional requirements describe what the system must do. They are split into pipeline-side requirements (the offline training subsystem) and demo-side requirements (the web application). The first demo requirement is the core executable functionality, per-flow classification of submitted traffic, whose end-to-end operation must be demonstrated.

\subsection{Pipeline Requirements}
\label{subsec:pipeline_requirements}

\par Ingestion must be reproducible (FR-P1): the pipeline ingests the raw per-folder CICIoT2023 CSVs, attaches the folder name as the fine-grained class label, drops exact duplicate rows, and writes a single compact columnar intermediate, such that re-running with the same seed and source files yields a byte-identical intermediate. Sampling must be uniform per class (FR-P2): for each fine-grained label the pipeline randomly subsamples at most a configurable per-class cap of rows with a fixed seed, drawing uniformly rather than taking the leading rows, since the leading rows of each attack folder come from a single capture session and would bias the model toward that session. The pipeline must produce three splits over the same subsampled rows (FR-P3): a temporal split (source captures per folder sorted, earliest 70\% for training, next 15\% for validation, latest 15\% for test); a per-capture split (no source capture appears in more than one partition); and a random-row stratified split retained for parity with published numbers. For each split it must train and evaluate one model per granularity (FR-P4), reusing the same row indices across granularities and changing only the label encoding.

\par Imbalance is handled by class-weighted cross-entropy with weights from the training-set class frequencies (FR-P5); synthetic oversampling and weighted batch sampling are not used; Section~\ref{subsec:class_weighted_ce} gives the reasons. The pipeline must also train a Random Forest and a gradient-boosted-tree classifier on the same splits and report the same metrics (FR-P6), so the deep model can be compared against strong tabular baselines. For each (split, granularity) pair it reports accuracy, macro-F1, weighted-F1, macro-precision, macro-recall, the per-class precision/recall/F1 report, and the row-normalised confusion matrix (FR-P7). It measures single-flow inference latency at batch sizes 1, 32, 256, and 1024 over at least 10{,}000 warm runs, reporting p50, p95, and p99 end-to-end, with feature scaling included in the timed path (FR-P8); a mean latency alone does not satisfy this requirement. Finally, for each (split, granularity) pair it persists three artefacts, the fitted feature scaler, the fitted label encoder, and the trained model weights (FR-P9), in a format the serving subsystem can reload without retraining and under a naming convention uniform enough that new variants need no code changes (see Section~\ref{subsec:high_level_architecture}).

\subsection{Demo Requirements}
\label{subsec:demo_requirements}

\par The core demo requirement is per-flow classification of submitted traffic: the demo accepts either a CSV already in the CIC feature format or a raw packet capture, and returns a per-flow classification table with an aggregated summary. For packet-capture input the established CIC flow-extraction tool is invoked as a subprocess to derive the features; flow extraction is not reimplemented in Python, because that would risk subtle numerical drift between the live and training distributions. The remaining demo requirements support this core. Model-variant selection (FR-D2) presents a dropdown of every (split, granularity) pair with trained artefacts on disk and lets the user switch variant without restarting. The output schema is deterministic (FR-D3): a one-line summary giving the total flow count and a breakdown by predicted label sorted by descending count, followed by a table whose rows carry the input row index, the predicted label, and the maximum-softmax confidence to four decimal digits. Input is self-validating (FR-D4): an uploaded CSV is checked for the expected feature columns and a missing-column error names them in plain language, while a flow-extraction failure surfaces the tool's own error message rather than crashing the request.

\par Raw traffic is never persisted (FR-D5): submitted captures and CSVs are processed in a temporary location and removed when the request ends, and payloads are not logged to disk. Derived result metadata (the model used, the predicted label, the confidence, the classification mode, and a timestamp) may be persisted for an authenticated user to support the history feature, since it contains no traffic payload. The application is deployed as a permanent public service (FR-D6) reachable at a fixed HTTPS URL with no presenter-side configuration: the backend runs as a Docker container on Hugging Face Spaces and the React frontend on Vercel, with CORS configured for cross-origin calls, while a local Gradio tab UI remains as a development and fallback interface. The deployed frontend requires authentication and keeps a per-user history of past classifications (FR-D7, metadata only per FR-D5), with each user seeing only their own history through isolation enforced at the storage layer rather than in client code.

\section{Non-Functional Requirements}
\label{sec:non_functional_requirements}

\par The non-functional requirements describe how well the system must perform. End-to-end per-flow latency through the demo, covering CSV parsing, feature scaling, the forward pass, and result formatting, must stay below 100~ms at p95 on a consumer laptop CPU for batches of up to 1024 flows (NFR-1); packet-capture parsing is excluded from this budget because the flow-extraction tool dominates that path, and capture-to-verdict latency is reported separately. For classification performance (NFR-2), the temporal-split model must reach at least 0.85 weighted-F1 on the held-out binary test set, which is the operational acceptance bar for the core detection capability; the harder 8-class attack-family task is a secondary objective whose weighted-F1 is expected to be lower, since the features describe transport-layer behaviour and cannot fully separate the application-layer families, and the temporal headline number is expected to sit 2 to 10 points below the random-split number, which is treated as the realistic generalisation story.

\par The pipeline must be reproducible (NFR-3): every random source (numerical libraries, the deep-learning framework on CPU and GPU, and the dataset samplers) is seeded from a single deterministic seed, so re-running produces identical splits and class weights and metrics within $\pm 0.001$ F1. It must fit within 16~GB of RAM at every stage (NFR-4), which rules out materialising the full multi-tens-of-millions-of-rows CSV corpus and forces lazy columnar processing with on-disk intermediates. It must run on CPU-only hardware (NFR-5), using GPU acceleration opportunistically but never as a hard dependency. The demo must run on Windows and Linux (NFR-6), the only platform-specific concern being the path of the flow-extraction launcher script, and the backend must be packageable as a Docker image for Hugging Face Spaces. Operationally it must stay simple (NFR-7): a permanent HTTPS URL with no presenter-side process management, a single command for local development, and no database, message queue, or dedicated model server, since the backend loads artefacts directly into the application process at startup. Reporting must be faithful (NFR-8): the headline metrics are taken from the temporal split, even though the random split would report higher numbers, and per-class metrics always accompany aggregate metrics so majority-class dominance cannot mask weak minority-class performance.

\section{Use Cases}
\label{sec:use_cases}

\par The principal interactions with the system are summarised in Table~\ref{tab:use_cases}. The core classification path, handling a captured packet trace, is then narrated in full as the representative flow; the others follow the same request shape.

\begin{table}[htbp]
\begin{center}
\begin{tabular}{|p{34pt}|p{92pt}|p{95pt}|p{125pt}|}
\hline
\textbf{ID} & \textbf{Actor} & \textbf{Trigger} & \textbf{Outcome} \\
\hline
\hline UC-1 & Researcher & Runs the training pipeline & A complete artefact triple (scaler, encoder, weights) per requested (split, granularity) pair, ready for serving. \\
\hline UC-2 & Operator & Uploads a packet capture & Per-flow verdict table and summary; temporary capture removed afterwards. \\
\hline UC-3 & Operator & Uploads a pre-extracted CIC CSV & Per-flow verdict table; missing columns reported as a readable error. \\
\hline UC-4 & Operator & Selects a different model card & Result screen re-rendered under the new variant with no restart or reload. \\
\hline UC-5 & Operator & Opens the public NEURAL IDS URL & Demo reachable from any network at a permanent URL with no presenter-side setup. \\
\hline
\end{tabular}
\end{center}
\caption{Principal use cases with their actors, triggers, and outcomes.}
\label{tab:use_cases}
\end{table}

\par The core classification flow (UC-2) runs as follows. The operator opens the public NEURAL IDS URL, selects a model type, switches to the packet-capture tab, and uploads a capture produced by the attack-script step. The application copies the capture into a temporary location, invokes the flow-extraction tool as a subprocess to obtain a feature CSV, applies the scaler fitted on the training set, runs the forward pass, applies softmax, and inverse-transforms the predicted indices to label strings, then renders the summary line and the per-flow verdict table. The temporary location is removed afterwards and nothing is persisted to disk. If the flow-extraction tool exits non-zero, the application returns its captured stderr with an empty table; if no flow-extraction backend is available, it returns a configuration message explaining how to provide one. The pre-extracted-CSV path (UC-3) is identical except that it validates the CIC column set instead of running flow extraction.

\section{System Specification}
\label{sec:system_specification}

\par The system specification translates the requirements above into the high-level architecture, the technology choices, and the interfaces that govern how components cooperate. The detailed design and implementation are presented in Chapter~\ref{chap:design_implementation}; this section fixes the contract.

\subsection{High-Level Architecture}
\label{subsec:high_level_architecture}

\par The system consists of two pipelines that meet at a small set of serialised artefacts. The training pipeline runs offline and produces those artefacts. The serving pipeline is the live demo, which loads them and answers classification requests. Figure~\ref{fig:system_architecture} shows the data flow.

\begin{figure}[htbp]
    \centering
    \resizebox{\textwidth}{!}{%
    \begin{tikzpicture}[
        node distance=0.6cm and 0.8cm,
        layer/.style={
            rectangle,
            rounded corners=4pt,
            draw=black!70,
            thick,
            minimum width=2.6cm,
            minimum height=1.0cm,
            text width=2.5cm,
            align=center,
            font=\small
        },
        train/.style={layer, fill=blue!8},
        artefact/.style={layer, fill=yellow!18, minimum width=2.4cm},
        serve/.style={layer, fill=green!10},
        arrow/.style={
            ->,
            >=latex,
            thick,
            black!60
        }
    ]
    % Training pipeline (top row)
    \node[train] (raw)        {Raw CIC\\captures};
    \node[train, right=of raw]        (parquet)   {Cleaned\\columnar\\intermediate};
    \node[train, right=of parquet]    (split)     {Splits\\(temporal /\\per-capture /\\random)};
    \node[train, right=of split]      (train)     {Model\\training};

    % Artefact layer (middle)
    \node[artefact, below=of train, yshift=-0.2cm] (scaler)  {Fitted\\scaler};
    \node[artefact, left=of scaler]                (encoder) {Fitted\\encoder};
    \node[artefact, left=of encoder]               (weights) {Trained\\weights};

    % Serving pipeline (bottom row)
    \node[serve, below=of weights, yshift=-0.2cm]  (pcap)    {Capture /\\CSV upload};
    \node[serve, right=of pcap]                    (cic)     {Flow\\extractor};
    \node[serve, right=of cic]                     (predict) {Inference\\component};
    \node[serve, right=of predict]                 (gradio)  {Web UI};

    % Arrows: training row
    \draw[arrow] (raw) -- (parquet);
    \draw[arrow] (parquet) -- (split);
    \draw[arrow] (split) -- (train);

    % Train -> artefacts
    \draw[arrow] (train.south) -- (scaler.north);
    \draw[arrow] (train.south) -- (encoder.north);
    \draw[arrow] (train.south) -- (weights.north);

    % Artefacts -> predictor
    \draw[arrow] (scaler.south) -- (predict.north);
    \draw[arrow] (encoder.south) -- (predict.north);
    \draw[arrow] (weights.south) -- (predict.north);

    % Serving row
    \draw[arrow] (pcap) -- (cic);
    \draw[arrow] (cic) -- (predict);
    \draw[arrow] (predict) -- (gradio);

    \end{tikzpicture}%
    }
    \caption{High-level architecture of the RT-IDS application. The training pipeline (blue) produces three serialised artefacts (yellow), which the serving pipeline (green) loads to answer per-flow classification requests.}
    \label{fig:system_architecture}
\end{figure}

\subsection{Technology Stack and Rationale}
\label{subsec:tech_stack}

\par The technology choices are driven by the requirements above. Table~\ref{tab:tech_stack} lists each concern, the selected technology, and the reason it was chosen over the obvious alternatives.

\begin{table}[htbp]
\begin{center}
\small
\begin{tabular}{|p{70pt}|p{82pt}|p{198pt}|}
\hline
\textbf{Concern} & \textbf{Choice} & \textbf{Rationale} \\
\hline
\hline
Tabular data wrangling & Lazy columnar engine, compressed columnar intermediates & Lazy scanning lets the multi-tens-of-millions-of-rows CSV corpus be deduplicated and subsampled without ever materialising it in memory; columnar storage with strong compression cuts disk usage and reload time. \\
\hline
Feature scaling & Robust scaler (median / IQR) & CIC features have heavy-tailed distributions (rate, IAT, total bytes); scaling by median and IQR is more stable than zero-mean / unit-variance scaling under such tails. \\
\hline
Label encoding & Integer label encoder & Compact and reversible; the inverse transform is needed at inference time to display human-readable labels. \\
\hline
Imbalance handling & Class-weighted cross-entropy & Combats severe imbalance without synthetic data; SMOTE and weighted batch sampling are rejected (Section~\ref{subsec:class_weighted_ce}). \\
\hline
Deep model & Two-hidden-layer MLP with batch normalisation and dropout & MLPs are the natural fit for fixed-length tabular inputs, train fast on CPU, and serve as the strong baseline expected by the IDS literature. \\
\hline
Tree baselines & Random Forest, gradient-boosted trees & Required to demonstrate that the deep model is not strictly worse than a tree, which is the first question committees ask about tabular tasks. \\
\hline
Capture $\to$ flow CSV & CICFlowMeter (the CIC team's reference tool) & Reimplementing flow aggregation in Python would risk subtle numerical drift between the live-demo features and the training distribution. \\
\hline
Demo UI & Gradio UI + FastAPI REST endpoint; React dashboard (NEURAL IDS) & Gradio serves the local tab interface; FastAPI exposes a REST endpoint consumed by the React frontend, which adds model cards, a confidence gauge, and session history. \\
\hline
Public exposure & Docker on Hugging Face Spaces; React on Vercel & Backend containerised and served via HF Spaces; React frontend on Vercel. No router configuration required; permanent HTTPS URL. \\
\hline
Reproducibility & Single deterministic seed propagated to every random source & Required by NFR-3. \\
\hline
\end{tabular}
\end{center}
\caption{Technology stack of the RT-IDS application and rationale for each choice.}
\label{tab:tech_stack}
\end{table}

\subsection{Constraints and Assumptions}
\label{subsec:constraints}

\par Several constraints and assumptions bound the work and are acknowledged in the methodology. The 39-feature public CSV release of CICIoT2023 is the only source of training data; the 46-feature variant used by some published comparisons was not available from the official UNB CIC distribution at the time of writing, and although it is mirrored on third-party platforms, its provenance and preprocessing differ, so numerical results are not directly comparable across the two feature sets. The dataset was captured in a smart-home testbed of 105 IoT devices, so generalisation to industrial control systems, medical IoT, or vehicular networks is not claimed. Its features come from headers and metadata only, so detection of heavily encrypted traffic and of application-layer web attacks (SQL injection, XSS, command injection) is expected to be weak; this follows from the transport-layer feature set rather than from a model defect and is discussed in the results chapter. The production serving target is a single-container backend on Hugging Face Spaces with a static Vercel frontend, so distributed serving, load balancing, and high availability are out of scope and the container's cold-start latency is tolerated. Authentication uses Supabase email/password to restrict access during the thesis period and to scope each user's history (FR-D7), limited to identity verification with no role-based permissions or audit logging, and per-user isolation is enforced by Postgres row-level security rather than in client code, while the local-development Gradio UI does not authenticate. Finally, the deployed model is the one trained offline by the pipeline; the demo does not adapt to feedback during a session.
